
% Copyright Glen Newton glen.newton@gmail.com
%
\documentclass[12pt]{article}
\usepackage{imakeidx}
\usepackage{hyperref}
\usepackage[hyphenbreaks]{breakurl}
\usepackage[most]{tcolorbox}
%\usepackage{savetrees}
\usepackage{graphicx}
\usepackage[utf8]{inputenc}
\usepackage[english]{babel}
\usepackage{multicol}
\usepackage{seqsplit}
\usepackage{xcolor}
\usepackage{etoolbox}
\usepackage{tocloft}
\usepackage{fancyhdr}
\usepackage{listings}
\usepackage{tabularx}
\usepackage[export]{adjustbox}
\usepackage{multicol}

\usepackage[inner=10mm,outer=10mm]{geometry}


%\usepackage{../sty/azureicons} 

\graphicspath{{../icons_tex/}}

%%%%%%%%%%%%%%%%%%%%%%%%%%%
\title{\LaTeX\ Style for Azure Icons}
\author{Glen Newton \\\texttt{\href{mailto:glen.newton@gmail.com}{glen.newton@gmail.com}}}

\date{\today}

%%%%%%%%%%%%%%%%%%%%%%%%%%%%%%%%%%%%%%%%%%%%%%%%%%%%%%%%%%%%%%%%%%%%%%%%
%% Change below lines to change look of document %%%%%%%%%%%%%%%%%%%%%%%
%%%%%%%%%%%%%%%%%%%%%%%%%%%%%%%%%%%%%%%%%%%%%%%%%%%%%%%%%%%%%%%%%%%%%%%%

\pagestyle{fancy}
\fancyhf{}
\rfoot{\hyperref[index]{General Index}}
\cfoot{\hyperref[macros]{Macro Index}}
\lfoot{\hyperref[toc]{TOC}}

\rhead{Page \thepage}
\chead{\href{https://github.com/gnewton/azureArchIcons2LaTeX}{azureArchIcons2LaTeX}}

\newcommand{\iconsize}{1cm}
\newcommand{\entrySpacing}{2mm}
\newcommand{\minipagespacing}{8cm}
%\renewcommand{\familydefault}{\sfdefault}

%% Arch and Res sections start/end
\newcommand{\archStart}{\section{Services}\begin{multicols}{2}\footnotesize}
\newcommand{\archEnd}{\end{multicols}\clearpage}
\newcommand{\resStart}{\newpage\section{Resources}\begin{multicols}{2}\footnotesize}
\newcommand{\resEnd}{\end{multicols}\clearpage}

%% Command used for each icon entry
%\newcommand{\gxs}[5]{\vcenteredinclude{#2}{\hspace{3mm} \begin{minipage}{\minipagespacing}{\bf \small{#1}}\\ {\footnotesize \tt \seqsplit{#3}} \\\tt Macro:\seqsplit{#4}\index[macros]{#4 (#5)} \end{minipage}}\vspace{\vspacing}\newline}

% 1 English name
% 2 graphic
% 3 PDF file (printed TeX)
% 4 Macro name
% 5 Group
\newcommand{\gxs}[5]{
\noindent\begin{tabularx}{9cm}{p{12mm}p{7.8cm}}
               {#2}& 
               #1\index{#1 (#5)}\newline
               \texttt{#3}\newline
               Macro: {\textbackslash}#4 \index[macros]{#4 (#5)}\newline \\
\end{tabularx}
\vspace{\entrySpacing}
}

%% From: https://tex.stackexchange.com/a/847
\hypersetup{
    colorlinks,
    linkcolor={red!50!black},
    citecolor={blue!50!black},
    urlcolor={blue!80!black}
}

%%%%%%%%%%%%%%%%%%%%%%%%%%%%%%%%%%%%%%%%%%%%%%%%%%%%%%%%%%%%%%%%%%%%%%%%
%% Do not change rest of preamble (below this line) unless you know what
%% you are doing %%%%%%%%%%%%%%%%%%%%%%%%%%%%%%%%%%%%%%%%%%%%%%%%%%%%%%%
%%%%%%%%%%%%%%%%%%%%%%%%%%%%%%%%%%%%%%%%%%%%%%%%%%%%%%%%%%%%%%%%%%%%%%%%

%% From: https://codepunk.io/change-the-background-color-of-a-blockquote-in-latex-2/
\definecolor{block-gray}{gray}{.75}
\newtcolorbox{blockquote}{colback=block-gray,grow to right by=-1mm,grow to left by=-1mm,boxrule=0pt,boxsep=0pt,breakable}

\renewcommand{\cftsecleader}{\cftdotfill{\cftdotsep}}

\makeatletter
\appto{\endmulticols}{\@doendpe}
\makeatother

%% Ref to index: From https://tex.stackexchange.com/questions/182710/how-to-reference-the-index
\makeindex[intoc,title={Index: General\label{index}}]
\makeindex[name=macros, intoc,title={Index: \LaTeX \ \ Macros\label{macros}}]

%% "Too many open files" problem: From: https://stackoverflow.com/questions/1715677/error-too-many-open-files-in-pdflatex/1720556#1720556
\let\mypdfximage\pdfximage
\def\pdfximage{\immediate\mypdfximage}

 \newcommand{\vcenteredinclude}[1]{\begingroup
   \setbox0=\hbox{#1}%
   \parbox[t]{\wd0}{\box0}\endgroup}

%%%%%%%%%%%%%%%%%%%%%%%%%%%%%%%%%%%%%%%%%%%%%%%%%%%%%%%%%%%%%%%%%%%%%%%
\begin{document}

\maketitle
\tableofcontents\label{toc}
%\clearpage
\noindent\makebox[\linewidth]{\rule{\paperwidth}{0.4pt}}

%%%%%%
\section{What is this?}
\LaTeX\ style (\texttt{awsicons.sty}) which allows you to use the \href{https://aws.amazon.com/architecture/icons/}{AWS architectural icons} in your \LaTeX\  documents.

\vspace{3mm}

%\noindent The AWS SVG icons are converted to {\LaTeX}--compatible PDF form using \href{https://inkscape.org/}{Inkscape} \texttt{\inkscapeVersion} using the purpose--written Go program \href{https://github.com/gnewton/awsArchIcons2LaTeX/aws2tex}{aws2tex}. 

\noindent This Go program, \texttt{aws2tex}, creates:
\begin{itemize}
\item \texttt{awsicons.sty} file
\item \LaTeX--compatible PDF versions of icons (and written to the \texttt{icons\_tex} directory, using Inkscape).
\item This document, which is to illustrate as well as index these icons and their corresponding \LaTeX\ \  macros.
  \newline
  \textbf{NB:} Not all icons are converted:
  \begin{itemize}
  \item For \textit{Resources}, only the \texttt{\_Light} versions. The \texttt{\_Dark} are omitted.
  \item For \textit{Services}, only the \texttt{\_64} versions. The lower resolution \texttt{\_16,\_32,\_48} are omitted.
  \end{itemize}
For instructions on how to generate this document with \textit{all} icons, please see the instructions in the \href{https://github.com/gnewton/awsArchIcons2LaTeX}{README.md} at the github repo.\\
This document is: \url{https://github.com/gnewton/awsArchIcons2LaTeX/awsAllIcons.pdf}
\end{itemize}
This document also includes hyperlinks to Google searches of the titles of each icon.


%%%%%%
\section{How to use it?}
There are two parts:
\begin{enumerate}
\item \texttt{awsicons.sty} file
\item directory with all of the icons rendered into \LaTeX\ compatible form (in the repo, this is: \texttt{icons\_tex}).
\end{enumerate}
All you need to do is use the package to your \LaTeX\ document using \texttt{{\textbackslash}usepackage} (after having added the \texttt{awsicons.sty} file to where \LaTeX\ can see it, or put it in the same directory as your document) AND add the directory containing all the generated icons to the \texttt{{\textbackslash}graphicspath\{...\}}.\\

\noindent For example:
\vspace{4mm}

\begin{blockquote}
\begin{verbatim}
% preamble
\usepackage{awsicons} 
\graphicspath{../icons_tex/}
% after \begin{document}
Here is a 1cm API Gateway:
\RAPIGatewayEndpoint{1cm}
Here is a 3cm Lambda function:
\RLambdaLambdaFunc{3cm}

\end{verbatim}
\end{blockquote}

\noindent Resulting in:\\
\begin{blockquote}
Here is a 1cm API Gateway:
%\RAPIGatewayEndpoint{1cm}
Here is a 3cm Lambda function:
%\RLambdaLambdaFunc{3cm}

\end{blockquote}

[Without the gray background]
    
%%%%%%
\subsection{Examples}
There are more complex examples in the \href{https://github.com/gnewton/awsArchIcons2LaTeX/tree/main/examples}{examples} directory, including one which recreates an AWS Architecture blog diagram.

\begin{itemize}
\item \href{https://github.com/gnewton/awsArchIcons2LaTeX/blob/main/examples/simple.pdf}{Simple example}
\item AWS Architecture blog diagram: \textit{\href{https://github.com/gnewton/awsArchIcons2LaTeX/blob/main/examples/Data-pipeline-Grov-Technologies.pdf}{Building a Controlled Environment Agriculture Platform}\footnote{\url{https://aws.amazon.com/blogs/architecture/building-a-controlled-environment-agriculture-platform/}}}, by Ashu Joshi, 2020.12.22.
\end{itemize}
  
\section{Sources}
The AWS Architecture icons used here are from \url{https://aws.amazon.com/architecture/icons/}.\\
%
%The specific assets package used: \href{https://d1.awsstatic.com/webteam/architecture-icons/Q32020/\assetZipFile}{\assetZipFileSplit}
%The specific assets package used: {\sloppy \href{https://d1.awsstatic.com/webteam/architecture-icons/Q32020/}{\assetZipFile}}

\vspace{3cm}
\noindent \copyright\ Glen Newton 2020\\
Icons in this document (I believe) are: \ \copyright\ Amazon Web Services (AWS)
\vspace{5mm}

\setlength{\parindent}{0pt}

\begin{multicols}{2}\footnotesize
%%%%%%%%%%%%% ICON CONTENT: generated by azure2tex:
\subsection{DevOps}
\gxs{Application Insights}{\includegraphics[width=1cm,valign=t]{00012_ApplicationInsights.pdf}}{00012\_ApplicationInsights.pdf}{AppInsights}{DevOps}

\gxs{Azure DevOps}{\includegraphics[width=1cm,valign=t]{10261_AzureDevOps.pdf}}{10261\_AzureDevOps.pdf}{AzrDevOps}{DevOps}

\gxs{DevTest Labs}{\includegraphics[width=1cm,valign=t]{10264_DevTestLabs.pdf}}{10264\_DevTestLabs.pdf}{DevTestLabs}{DevOps}

\gxs{Lab Services}{\includegraphics[width=1cm,valign=t]{10265_LabServices.pdf}}{10265\_LabServices.pdf}{LabSvcs}{DevOps}

\subsection{Integration}
\gxs{SQL Data Warehouses}{\includegraphics[width=1cm,valign=t]{00036_SQLDataWarehouses.pdf}}{00036\_SQLDataWarehouses.pdf}{SQLDataWarehouses}{Integration}

\gxs{API Management Services}{\includegraphics[width=1cm,valign=t]{10042_APIManagementServices.pdf}}{10042\_APIManagementServices.pdf}{APIMngmtSvcs}{Integration}

\gxs{Logic Apps}{\includegraphics[width=1cm,valign=t]{10201_LogicApps.pdf}}{10201\_LogicApps.pdf}{LogicApps}{Integration}

\gxs{Event Grid Topics}{\includegraphics[width=1cm,valign=t]{10206_EventGridTopics.pdf}}{10206\_EventGridTopics.pdf}{EventGridTopics}{Integration}

\gxs{Relays}{\includegraphics[width=1cm,valign=t]{10209_Relays.pdf}}{10209\_Relays.pdf}{Relays}{Integration}

\gxs{Software as a Service}{\includegraphics[width=1cm,valign=t]{10213_SoftwareasaService.pdf}}{10213\_SoftwareasaService.pdf}{SoftwareasaSvc}{Integration}

\gxs{Event Grid Domains}{\includegraphics[width=1cm,valign=t]{10215_EventGridDomains.pdf}}{10215\_EventGridDomains.pdf}{EventGridDomains}{Integration}

\gxs{Azure Data Catalog}{\includegraphics[width=1cm,valign=t]{10216_AzureDataCatalog.pdf}}{10216\_AzureDataCatalog.pdf}{AzrDataCatalog}{Integration}

\gxs{Integration Accounts}{\includegraphics[width=1cm,valign=t]{10218_IntegrationAccounts.pdf}}{10218\_IntegrationAccounts.pdf}{IntegAccnts}{Integration}

\gxs{Event Grid Subscriptions}{\includegraphics[width=1cm,valign=t]{10221_EventGridSubscriptions.pdf}}{10221\_EventGridSubscriptions.pdf}{EventGridSubscriptions}{Integration}

\subsection{Intune}
\gxs{Device Security Apple}{\includegraphics[width=1cm,valign=t]{00399_DeviceSecurityApple.pdf}}{00399\_DeviceSecurityApple.pdf}{DevSecApple}{Intune}

\gxs{Device Security Google}{\includegraphics[width=1cm,valign=t]{00399_DeviceSecurityGoogle.pdf}}{00399\_DeviceSecurityGoogle.pdf}{DevSecGoogle}{Intune}

\gxs{Device Security Windows}{\includegraphics[width=1cm,valign=t]{00399_DeviceSecurityWindows.pdf}}{00399\_DeviceSecurityWindows.pdf}{DevSecWin}{Intune}

\gxs{Azure AD Roles and Administrators}{\includegraphics[width=1cm,valign=t]{10340_AzureADRolesandAdministrators.pdf}}{10340\_AzureADRolesandAdministrators.pdf}{AzrADRolesandAdministrators}{Intune}

\subsection{Web}
\gxs{Notification Hub Namespaces}{\includegraphics[width=1cm,valign=t]{10053_NotificationHubNamespaces.pdf}}{10053\_NotificationHubNamespaces.pdf}{NotifHubNS}{Web}

\gxs{Azure Media Service}{\includegraphics[width=1cm,valign=t]{10309_AzureMediaService.pdf}}{10309\_AzureMediaService.pdf}{AzrMediaSvc}{Web}

\subsection{Analytics}
\gxs{Log Analytics Workspaces}{\includegraphics[width=1cm,valign=t]{00009_LogAnalyticsWorkspaces.pdf}}{00009\_LogAnalyticsWorkspaces.pdf}{LogAnalyticsWorkspaces}{Analytics}

\gxs{Event Hubs}{\includegraphics[width=1cm,valign=t]{00039_EventHubs.pdf}}{00039\_EventHubs.pdf}{EventHubs}{Analytics}

\gxs{Stream Analytics Jobs}{\includegraphics[width=1cm,valign=t]{00042_StreamAnalyticsJobs.pdf}}{00042\_StreamAnalyticsJobs.pdf}{StrmAnalyticsJobs}{Analytics}

\gxs{Azure Synapse Analytics}{\includegraphics[width=1cm,valign=t]{00606_AzureSynapseAnalytics.pdf}}{00606\_AzureSynapseAnalytics.pdf}{AzrSynapseAnalytics}{Analytics}

\gxs{HD Insight Clusters}{\includegraphics[width=1cm,valign=t]{10142_HDInsightClusters.pdf}}{10142\_HDInsightClusters.pdf}{HDInsightClusters}{Analytics}

\gxs{Analysis Services}{\includegraphics[width=1cm,valign=t]{10148_AnalysisServices.pdf}}{10148\_AnalysisServices.pdf}{AnalysisSvcs}{Analytics}

\gxs{Event Hub Clusters}{\includegraphics[width=1cm,valign=t]{10149_EventHubClusters.pdf}}{10149\_EventHubClusters.pdf}{EventHubClusters}{Analytics}

\gxs{Data Lake Store Gen1}{\includegraphics[width=1cm,valign=t]{10150_DataLakeStoreGen1.pdf}}{10150\_DataLakeStoreGen1.pdf}{DataLakeStoreGenOne}{Analytics}

\subsection{Blockchain}
\gxs{Outbound Connection}{\includegraphics[width=1cm,valign=t]{10364_OutboundConnection.pdf}}{10364\_OutboundConnection.pdf}{OutboundConnection}{Blockchain}

\gxs{Azure Blockchain Service}{\includegraphics[width=1cm,valign=t]{10366_AzureBlockchainService.pdf}}{10366\_AzureBlockchainService.pdf}{AzrBlockchainSvc}{Blockchain}

\gxs{Azure Token Service}{\includegraphics[width=1cm,valign=t]{10367_AzureTokenService.pdf}}{10367\_AzureTokenService.pdf}{AzrTokenSvc}{Blockchain}

\gxs{ABS Member}{\includegraphics[width=1cm,valign=t]{10374_ABSMember.pdf}}{10374\_ABSMember.pdf}{ABSMember}{Blockchain}

\gxs{Consortium}{\includegraphics[width=1cm,valign=t]{10375_Consortium.pdf}}{10375\_Consortium.pdf}{Consortium}{Blockchain}

\subsection{Compute}
\gxs{Disk Encryption Sets}{\includegraphics[width=1cm,valign=t]{00398_DiskEncryptionSets.pdf}}{00398\_DiskEncryptionSets.pdf}{DiskEncrSets}{Compute}

\gxs{Workspaces}{\includegraphics[width=1cm,valign=t]{00400_Workspaces.pdf}}{00400\_Workspaces.pdf}{Workspaces}{Compute}

\gxs{Virtual Machine}{\includegraphics[width=1cm,valign=t]{10021_VirtualMachine.pdf}}{10021\_VirtualMachine.pdf}{VirtMachine}{Compute}

\gxs{Kubernetes Services}{\includegraphics[width=1cm,valign=t]{10023_KubernetesServices.pdf}}{10023\_KubernetesServices.pdf}{KuberSvcs}{Compute}

\gxs{Mesh Applications}{\includegraphics[width=1cm,valign=t]{10024_MeshApplications.pdf}}{10024\_MeshApplications.pdf}{MeshApps}{Compute}

\gxs{Availability Sets}{\includegraphics[width=1cm,valign=t]{10025_AvailabilitySets.pdf}}{10025\_AvailabilitySets.pdf}{AvailabilitySets}{Compute}

\gxs{Disks Snapshots}{\includegraphics[width=1cm,valign=t]{10026_DisksSnapshots.pdf}}{10026\_DisksSnapshots.pdf}{DisksSnapshots}{Compute}

\gxs{OS Images Classic}{\includegraphics[width=1cm,valign=t]{10027_OSImagesClassic.pdf}}{10027\_OSImagesClassic.pdf}{OSImgsClassic}{Compute}

\gxs{Virtual Machines Classic}{\includegraphics[width=1cm,valign=t]{10028_VirtualMachinesClassic.pdf}}{10028\_VirtualMachinesClassic.pdf}{VirtMachinesClassic}{Compute}

\gxs{Function Apps}{\includegraphics[width=1cm,valign=t]{10029_FunctionApps.pdf}}{10029\_FunctionApps.pdf}{FuncApps}{Compute}

\gxs{Cloud Services Classic}{\includegraphics[width=1cm,valign=t]{10030_CloudServicesClassic.pdf}}{10030\_CloudServicesClassic.pdf}{CloudSvcsClassic}{Compute}

\gxs{Batch Accounts}{\includegraphics[width=1cm,valign=t]{10031_BatchAccounts.pdf}}{10031\_BatchAccounts.pdf}{BatchAccnts}{Compute}

\gxs{Disks}{\includegraphics[width=1cm,valign=t]{10032_Disks.pdf}}{10032\_Disks.pdf}{Disks}{Compute}

\gxs{Images}{\includegraphics[width=1cm,valign=t]{10033_Images.pdf}}{10033\_Images.pdf}{Imgs}{Compute}

\gxs{VM Scale Sets}{\includegraphics[width=1cm,valign=t]{10034_VMScaleSets.pdf}}{10034\_VMScaleSets.pdf}{VMScaleSets}{Compute}

\gxs{App Services}{\includegraphics[width=1cm,valign=t]{10035_AppServices.pdf}}{10035\_AppServices.pdf}{AppSvcs}{Compute}

\gxs{Service Fabric Clusters}{\includegraphics[width=1cm,valign=t]{10036_ServiceFabricClusters.pdf}}{10036\_ServiceFabricClusters.pdf}{SvcFabricClusters}{Compute}

\gxs{Image Definitions}{\includegraphics[width=1cm,valign=t]{10037_ImageDefinitions.pdf}}{10037\_ImageDefinitions.pdf}{ImgDefinitions}{Compute}

\gxs{Image Versions}{\includegraphics[width=1cm,valign=t]{10038_ImageVersions.pdf}}{10038\_ImageVersions.pdf}{ImgVersions}{Compute}

\gxs{Shared Image Galleries}{\includegraphics[width=1cm,valign=t]{10039_SharedImageGalleries.pdf}}{10039\_SharedImageGalleries.pdf}{SharedImgGalleries}{Compute}

\gxs{VM Images Classic}{\includegraphics[width=1cm,valign=t]{10040_VMImagesClassic.pdf}}{10040\_VMImagesClassic.pdf}{VMImgsClassic}{Compute}

\gxs{Disks Classic}{\includegraphics[width=1cm,valign=t]{10041_DisksClassic.pdf}}{10041\_DisksClassic.pdf}{DisksClassic}{Compute}

\gxs{Container Services Deprecated}{\includegraphics[width=1cm,valign=t]{10049_ContainerServicesDeprecated.pdf}}{10049\_ContainerServicesDeprecated.pdf}{ContainerSvcsDeprecated}{Compute}

\gxs{Container Instances}{\includegraphics[width=1cm,valign=t]{10104_ContainerInstances.pdf}}{10104\_ContainerInstances.pdf}{ContainerInsts}{Compute}

\subsection{General}
\gxs{All Resources}{\includegraphics[width=1cm,valign=t]{10001_AllResources.pdf}}{10001\_AllResources.pdf}{AllRess}{General}

\gxs{Subscriptions}{\includegraphics[width=1cm,valign=t]{10002_Subscriptions.pdf}}{10002\_Subscriptions.pdf}{Subscriptions}{General}

\gxs{Reservations}{\includegraphics[width=1cm,valign=t]{10003_Reservations.pdf}}{10003\_Reservations.pdf}{Reservations}{General}

\gxs{Service Health}{\includegraphics[width=1cm,valign=t]{10004_ServiceHealth.pdf}}{10004\_ServiceHealth.pdf}{SvcHealth}{General}

\gxs{Information}{\includegraphics[width=1cm,valign=t]{10005_Information.pdf}}{10005\_Information.pdf}{Information}{General}

\gxs{Recent}{\includegraphics[width=1cm,valign=t]{10006_Recent.pdf}}{10006\_Recent.pdf}{Recent}{General}

\gxs{Resource Groups}{\includegraphics[width=1cm,valign=t]{10007_ResourceGroups.pdf}}{10007\_ResourceGroups.pdf}{ResGrp}{General}

\gxs{Marketplace}{\includegraphics[width=1cm,valign=t]{10008_Marketplace.pdf}}{10008\_Marketplace.pdf}{Marketplace}{General}

\gxs{Quickstart Center}{\includegraphics[width=1cm,valign=t]{10010_QuickstartCenter.pdf}}{10010\_QuickstartCenter.pdf}{QuickstartCenter}{General}

\gxs{Management Groups}{\includegraphics[width=1cm,valign=t]{10011_ManagementGroups.pdf}}{10011\_ManagementGroups.pdf}{MngmtGrp}{General}

\gxs{Help and Support}{\includegraphics[width=1cm,valign=t]{10013_HelpandSupport.pdf}}{10013\_HelpandSupport.pdf}{HelpandSupport}{General}

\gxs{Dashboard}{\includegraphics[width=1cm,valign=t]{10015_Dashboard.pdf}}{10015\_Dashboard.pdf}{Dashboard}{General}

\gxs{Cost Management}{\includegraphics[width=1cm,valign=t]{10019_CostManagement.pdf}}{10019\_CostManagement.pdf}{CostMngmt}{General}

\gxs{Biz Talk}{\includegraphics[width=1cm,valign=t]{10779_BizTalk.pdf}}{10779\_BizTalk.pdf}{BizTalk}{General}

\gxs{Blob Block}{\includegraphics[width=1cm,valign=t]{10780_BlobBlock.pdf}}{10780\_BlobBlock.pdf}{BlobBlock}{General}

\gxs{Blob Page}{\includegraphics[width=1cm,valign=t]{10781_BlobPage.pdf}}{10781\_BlobPage.pdf}{BlobPg}{General}

\gxs{Branch}{\includegraphics[width=1cm,valign=t]{10782_Branch.pdf}}{10782\_Branch.pdf}{Branch}{General}

\gxs{Browser}{\includegraphics[width=1cm,valign=t]{10783_Browser.pdf}}{10783\_Browser.pdf}{Browser}{General}

\gxs{Bug}{\includegraphics[width=1cm,valign=t]{10784_Bug.pdf}}{10784\_Bug.pdf}{Bug}{General}

\gxs{Builds}{\includegraphics[width=1cm,valign=t]{10785_Builds.pdf}}{10785\_Builds.pdf}{Builds}{General}

\gxs{Cache}{\includegraphics[width=1cm,valign=t]{10786_Cache.pdf}}{10786\_Cache.pdf}{Cache}{General}

\gxs{Code}{\includegraphics[width=1cm,valign=t]{10787_Code.pdf}}{10787\_Code.pdf}{Code}{General}

\gxs{Commit}{\includegraphics[width=1cm,valign=t]{10788_Commit.pdf}}{10788\_Commit.pdf}{Commit}{General}

\gxs{Controls}{\includegraphics[width=1cm,valign=t]{10789_Controls.pdf}}{10789\_Controls.pdf}{Controls}{General}

\gxs{Controls Horizontal}{\includegraphics[width=1cm,valign=t]{10790_ControlsHorizontal.pdf}}{10790\_ControlsHorizontal.pdf}{ControlsHorizontal}{General}

\gxs{Cost Alerts}{\includegraphics[width=1cm,valign=t]{10791_CostAlerts.pdf}}{10791\_CostAlerts.pdf}{CostAlerts}{General}

\gxs{Cost Analysis}{\includegraphics[width=1cm,valign=t]{10792_CostAnalysis.pdf}}{10792\_CostAnalysis.pdf}{CostAnalysis}{General}

\gxs{Cost Budgets}{\includegraphics[width=1cm,valign=t]{10793_CostBudgets.pdf}}{10793\_CostBudgets.pdf}{CostBudgets}{General}

\gxs{Counter}{\includegraphics[width=1cm,valign=t]{10794_Counter.pdf}}{10794\_Counter.pdf}{Counter}{General}

\gxs{Cubes}{\includegraphics[width=1cm,valign=t]{10795_Cubes.pdf}}{10795\_Cubes.pdf}{Cubes}{General}

\gxs{Dev Console}{\includegraphics[width=1cm,valign=t]{10796_DevConsole.pdf}}{10796\_DevConsole.pdf}{DevConsole}{General}

\gxs{Download}{\includegraphics[width=1cm,valign=t]{10797_Download.pdf}}{10797\_Download.pdf}{Download}{General}

\gxs{Error}{\includegraphics[width=1cm,valign=t]{10798_Error.pdf}}{10798\_Error.pdf}{Error}{General}

\gxs{Extensions}{\includegraphics[width=1cm,valign=t]{10799_Extensions.pdf}}{10799\_Extensions.pdf}{Extensions}{General}

\gxs{File}{\includegraphics[width=1cm,valign=t]{10800_File.pdf}}{10800\_File.pdf}{File}{General}

\gxs{Files}{\includegraphics[width=1cm,valign=t]{10801_Files.pdf}}{10801\_Files.pdf}{Files}{General}

\gxs{Folder Blank}{\includegraphics[width=1cm,valign=t]{10802_FolderBlank.pdf}}{10802\_FolderBlank.pdf}{FolderBlank}{General}

\gxs{Folder Website}{\includegraphics[width=1cm,valign=t]{10803_FolderWebsite.pdf}}{10803\_FolderWebsite.pdf}{FolderWebsite}{General}

\gxs{FTP}{\includegraphics[width=1cm,valign=t]{10804_FTP.pdf}}{10804\_FTP.pdf}{FTP}{General}

\gxs{Gear}{\includegraphics[width=1cm,valign=t]{10805_Gear.pdf}}{10805\_Gear.pdf}{Gear}{General}

\gxs{Globe}{\includegraphics[width=1cm,valign=t]{10806_Globe.pdf}}{10806\_Globe.pdf}{Globe}{General}

\gxs{Globe Error}{\includegraphics[width=1cm,valign=t]{10807_GlobeError.pdf}}{10807\_GlobeError.pdf}{GlobeError}{General}

\gxs{Globe Success}{\includegraphics[width=1cm,valign=t]{10808_GlobeSuccess.pdf}}{10808\_GlobeSuccess.pdf}{GlobeSuccess}{General}

\gxs{Globe Warning}{\includegraphics[width=1cm,valign=t]{10809_GlobeWarning.pdf}}{10809\_GlobeWarning.pdf}{GlobeWarn}{General}

\gxs{Guide}{\includegraphics[width=1cm,valign=t]{10810_Guide.pdf}}{10810\_Guide.pdf}{Guide}{General}

\gxs{Heart}{\includegraphics[width=1cm,valign=t]{10811_Heart.pdf}}{10811\_Heart.pdf}{Heart}{General}

\gxs{Image}{\includegraphics[width=1cm,valign=t]{10812_Image.pdf}}{10812\_Image.pdf}{Img}{General}

\gxs{Input Output}{\includegraphics[width=1cm,valign=t]{10813_InputOutput.pdf}}{10813\_InputOutput.pdf}{InputOutput}{General}

\gxs{Journey Hub}{\includegraphics[width=1cm,valign=t]{10814_JourneyHub.pdf}}{10814\_JourneyHub.pdf}{JourneyHub}{General}

\gxs{Launch Portal}{\includegraphics[width=1cm,valign=t]{10815_LaunchPortal.pdf}}{10815\_LaunchPortal.pdf}{LaunchPortal}{General}

\gxs{Learn}{\includegraphics[width=1cm,valign=t]{10816_Learn.pdf}}{10816\_Learn.pdf}{Learn}{General}

\gxs{Load Test}{\includegraphics[width=1cm,valign=t]{10817_LoadTest.pdf}}{10817\_LoadTest.pdf}{LoadTest}{General}

\gxs{Location}{\includegraphics[width=1cm,valign=t]{10818_Location.pdf}}{10818\_Location.pdf}{Location}{General}

\gxs{Log Streaming}{\includegraphics[width=1cm,valign=t]{10819_LogStreaming.pdf}}{10819\_LogStreaming.pdf}{LogStrming}{General}

\gxs{Management Portal}{\includegraphics[width=1cm,valign=t]{10820_ManagementPortal.pdf}}{10820\_ManagementPortal.pdf}{MngmtPortal}{General}

\gxs{Media File}{\includegraphics[width=1cm,valign=t]{10821_MediaFile.pdf}}{10821\_MediaFile.pdf}{MediaFile}{General}

\gxs{Mobile}{\includegraphics[width=1cm,valign=t]{10822_Mobile.pdf}}{10822\_Mobile.pdf}{Mob}{General}

\gxs{Mobile Engagement}{\includegraphics[width=1cm,valign=t]{10823_MobileEngagement.pdf}}{10823\_MobileEngagement.pdf}{MobEngagement}{General}

\gxs{Power}{\includegraphics[width=1cm,valign=t]{10824_Power.pdf}}{10824\_Power.pdf}{Power}{General}

\gxs{Powershell}{\includegraphics[width=1cm,valign=t]{10825_Powershell.pdf}}{10825\_Powershell.pdf}{Powershell}{General}

\gxs{Power Up}{\includegraphics[width=1cm,valign=t]{10826_PowerUp.pdf}}{10826\_PowerUp.pdf}{PowerUp}{General}

\gxs{Preview}{\includegraphics[width=1cm,valign=t]{10827_Preview.pdf}}{10827\_Preview.pdf}{Preview}{General}

\gxs{Process Explorer}{\includegraphics[width=1cm,valign=t]{10828_ProcessExplorer.pdf}}{10828\_ProcessExplorer.pdf}{ProcessExplorer}{General}

\gxs{Production Ready Database}{\includegraphics[width=1cm,valign=t]{10829_ProductionReadyDatabase.pdf}}{10829\_ProductionReadyDatabase.pdf}{ProductionReadyDB}{General}

\gxs{Resource Group List}{\includegraphics[width=1cm,valign=t]{10830_ResourceGroupList.pdf}}{10830\_ResourceGroupList.pdf}{ResGroupList}{General}

\gxs{Resource Linked}{\includegraphics[width=1cm,valign=t]{10831_ResourceLinked.pdf}}{10831\_ResourceLinked.pdf}{ResLinked}{General}

\gxs{Scale}{\includegraphics[width=1cm,valign=t]{10832_Scale.pdf}}{10832\_Scale.pdf}{Scale}{General}

\gxs{Scheduler}{\includegraphics[width=1cm,valign=t]{10833_Scheduler.pdf}}{10833\_Scheduler.pdf}{Scheduler}{General}

\gxs{Search}{\includegraphics[width=1cm,valign=t]{10834_Search.pdf}}{10834\_Search.pdf}{Srch}{General}

\gxs{Server Farm}{\includegraphics[width=1cm,valign=t]{10835_ServerFarm.pdf}}{10835\_ServerFarm.pdf}{ServerFarm}{General}

\gxs{Service Bus}{\includegraphics[width=1cm,valign=t]{10836_ServiceBus.pdf}}{10836\_ServiceBus.pdf}{SvcBus}{General}

\gxs{SSD}{\includegraphics[width=1cm,valign=t]{10837_SSD.pdf}}{10837\_SSD.pdf}{SSD}{General}

\gxs{Storage Azure Files}{\includegraphics[width=1cm,valign=t]{10838_StorageAzureFiles.pdf}}{10838\_StorageAzureFiles.pdf}{StorAzrFiles}{General}

\gxs{Storage Container}{\includegraphics[width=1cm,valign=t]{10839_StorageContainer.pdf}}{10839\_StorageContainer.pdf}{StorContainer}{General}

\gxs{Storage Queue}{\includegraphics[width=1cm,valign=t]{10840_StorageQueue.pdf}}{10840\_StorageQueue.pdf}{StorQueue}{General}

\gxs{Table}{\includegraphics[width=1cm,valign=t]{10841_Table.pdf}}{10841\_Table.pdf}{Table}{General}

\gxs{Tags}{\includegraphics[width=1cm,valign=t]{10842_Tags.pdf}}{10842\_Tags.pdf}{Tags}{General}

\gxs{TFS VC Repository}{\includegraphics[width=1cm,valign=t]{10843_TFSVCRepository.pdf}}{10843\_TFSVCRepository.pdf}{TFSVCRepository}{General}

\gxs{Toolbox}{\includegraphics[width=1cm,valign=t]{10844_Toolbox.pdf}}{10844\_Toolbox.pdf}{Toolbox}{General}

\gxs{Versions}{\includegraphics[width=1cm,valign=t]{10845_Versions.pdf}}{10845\_Versions.pdf}{Versions}{General}

\gxs{Website Power}{\includegraphics[width=1cm,valign=t]{10847_WebsitePower.pdf}}{10847\_WebsitePower.pdf}{WebsitePower}{General}

\gxs{Website Staging}{\includegraphics[width=1cm,valign=t]{10848_WebsiteStaging.pdf}}{10848\_WebsiteStaging.pdf}{WebsiteStaging}{General}

\gxs{Web Slots}{\includegraphics[width=1cm,valign=t]{10849_WebSlots.pdf}}{10849\_WebSlots.pdf}{WebSlots}{General}

\gxs{Web Test}{\includegraphics[width=1cm,valign=t]{10850_WebTest.pdf}}{10850\_WebTest.pdf}{WebTest}{General}

\gxs{Workbooks}{\includegraphics[width=1cm,valign=t]{10851_Workbooks.pdf}}{10851\_Workbooks.pdf}{Workbooks}{General}

\gxs{Workflow}{\includegraphics[width=1cm,valign=t]{10852_Workflow.pdf}}{10852\_Workflow.pdf}{Workflow}{General}

\gxs{Backlog}{\includegraphics[width=1cm,valign=t]{10853_Backlog.pdf}}{10853\_Backlog.pdf}{Backlog}{General}

\gxs{Media}{\includegraphics[width=1cm,valign=t]{10854_Media.pdf}}{10854\_Media.pdf}{Media}{General}

\gxs{Module}{\includegraphics[width=1cm,valign=t]{10855_Module.pdf}}{10855\_Module.pdf}{Module}{General}

\gxs{Search Grid}{\includegraphics[width=1cm,valign=t]{10856_SearchGrid.pdf}}{10856\_SearchGrid.pdf}{SrchGrid}{General}

\subsection{Security}
\gxs{Conditional Access}{\includegraphics[width=1cm,valign=t]{10233_ConditionalAccess.pdf}}{10233\_ConditionalAccess.pdf}{ConditionalAccess}{Security}

\gxs{Security Center}{\includegraphics[width=1cm,valign=t]{10241_SecurityCenter.pdf}}{10241\_SecurityCenter.pdf}{SecCenter}{Security}

\gxs{Application Security Groups}{\includegraphics[width=1cm,valign=t]{10244_ApplicationSecurityGroups.pdf}}{10244\_ApplicationSecurityGroups.pdf}{AppSecGrp}{Security}

\gxs{Key Vaults}{\includegraphics[width=1cm,valign=t]{10245_KeyVaults.pdf}}{10245\_KeyVaults.pdf}{KeyVlts}{Security}

\gxs{Azure Sentinel}{\includegraphics[width=1cm,valign=t]{10248_AzureSentinel.pdf}}{10248\_AzureSentinel.pdf}{AzrSentinel}{Security}

\gxs{ExtendedSecurityUpdates}{\includegraphics[width=1cm,valign=t]{10572_ExtendedSecurityUpdates.pdf}}{10572\_ExtendedSecurityUpdates.pdf}{ExtendedSecUpdates}{Security}

\subsection{Management + Governance}
\gxs{Monitor}{\includegraphics[width=1cm,valign=t]{00001_Monitor.pdf}}{00001\_Monitor.pdf}{Monitor}{Management + Governance}

\gxs{Alerts}{\includegraphics[width=1cm,valign=t]{00002_Alerts.pdf}}{00002\_Alerts.pdf}{Alerts}{Management + Governance}

\gxs{Advisor}{\includegraphics[width=1cm,valign=t]{00003_Advisor.pdf}}{00003\_Advisor.pdf}{Advisor}{Management + Governance}

\gxs{Blueprints}{\includegraphics[width=1cm,valign=t]{00006_Blueprints.pdf}}{00006\_Blueprints.pdf}{Blueprints}{Management + Governance}

\gxs{Activity Log}{\includegraphics[width=1cm,valign=t]{00007_ActivityLog.pdf}}{00007\_ActivityLog.pdf}{ActivityLog}{Management + Governance}

\gxs{Diagnostics Settings}{\includegraphics[width=1cm,valign=t]{00008_DiagnosticsSettings.pdf}}{00008\_DiagnosticsSettings.pdf}{DiagnosticsSettings}{Management + Governance}

\gxs{Log Analytics Workspaces}{\includegraphics[width=1cm,valign=t]{00009_LogAnalyticsWorkspaces.pdf}}{00009\_LogAnalyticsWorkspaces.pdf}{LogAnalyticsWorkspaces}{Management + Governance}

\gxs{Compliance}{\includegraphics[width=1cm,valign=t]{00011_Compliance.pdf}}{00011\_Compliance.pdf}{Compliance}{Management + Governance}

\gxs{Application Insights}{\includegraphics[width=1cm,valign=t]{00012_ApplicationInsights.pdf}}{00012\_ApplicationInsights.pdf}{AppInsights}{Management + Governance}

\gxs{My Customers}{\includegraphics[width=1cm,valign=t]{00014_MyCustomers.pdf}}{00014\_MyCustomers.pdf}{MyCustomers}{Management + Governance}

\gxs{Recovery Services Vaults}{\includegraphics[width=1cm,valign=t]{00017_RecoveryServicesVaults.pdf}}{00017\_RecoveryServicesVaults.pdf}{RecovSvcsVlts}{Management + Governance}

\gxs{Metrics}{\includegraphics[width=1cm,valign=t]{00020_Metrics.pdf}}{00020\_Metrics.pdf}{Metrics}{Management + Governance}

\gxs{Solutions}{\includegraphics[width=1cm,valign=t]{00021_Solutions.pdf}}{00021\_Solutions.pdf}{Solutions}{Management + Governance}

\gxs{Automation Accounts}{\includegraphics[width=1cm,valign=t]{00022_AutomationAccounts.pdf}}{00022\_AutomationAccounts.pdf}{AutomationAccnts}{Management + Governance}

\gxs{Operation Log Classic}{\includegraphics[width=1cm,valign=t]{00024_OperationLogClassic.pdf}}{00024\_OperationLogClassic.pdf}{OperationLogClassic}{Management + Governance}

\gxs{Service Providers}{\includegraphics[width=1cm,valign=t]{00025_ServiceProviders.pdf}}{00025\_ServiceProviders.pdf}{SvcPrvdr}{Management + Governance}

\gxs{Education}{\includegraphics[width=1cm,valign=t]{00026_Education.pdf}}{00026\_Education.pdf}{Education}{Management + Governance}

\gxs{Azure Lighthouse}{\includegraphics[width=1cm,valign=t]{00471_AzureLighthouse.pdf}}{00471\_AzureLighthouse.pdf}{AzrLighthouse}{Management + Governance}

\gxs{Azure Arc}{\includegraphics[width=1cm,valign=t]{00756_AzureArc.pdf}}{00756\_AzureArc.pdf}{AzrArc}{Management + Governance}

\gxs{User Privacy}{\includegraphics[width=1cm,valign=t]{10303_UserPrivacy.pdf}}{10303\_UserPrivacy.pdf}{UserPrivacy}{Management + Governance}

\gxs{Managed Applications Center}{\includegraphics[width=1cm,valign=t]{10313_ManagedApplicationsCenter.pdf}}{10313\_ManagedApplicationsCenter.pdf}{MngdAppsCenter}{Management + Governance}

\gxs{Policy}{\includegraphics[width=1cm,valign=t]{10316_Policy.pdf}}{10316\_Policy.pdf}{Policy}{Management + Governance}

\gxs{Resource Graph Explorer}{\includegraphics[width=1cm,valign=t]{10318_ResourceGraphExplorer.pdf}}{10318\_ResourceGraphExplorer.pdf}{ResGraphExplorer}{Management + Governance}

\gxs{MachinesAzureArc}{\includegraphics[width=1cm,valign=t]{10450_MachinesAzureArc.pdf}}{10450\_MachinesAzureArc.pdf}{MachinesAzrArc}{Management + Governance}

\subsection{Mixed Reality}
\gxs{Remote Rendering}{\includegraphics[width=1cm,valign=t]{00698_RemoteRendering.pdf}}{00698\_RemoteRendering.pdf}{RemoteRend}{Mixed Reality}

\subsection{Monitor}
\gxs{SAP Azure Monitor}{\includegraphics[width=1cm,valign=t]{00438_SAPAzureMonitor.pdf}}{00438\_SAPAzureMonitor.pdf}{SAPAzrMonitor}{Monitor}

\subsection{Azure Stack}
\gxs{Multi Tenancy}{\includegraphics[width=1cm,valign=t]{00965_MultiTenancy.pdf}}{00965\_MultiTenancy.pdf}{MultiTenancy}{Azure Stack}

\gxs{Infrastructure Backup}{\includegraphics[width=1cm,valign=t]{10108_InfrastructureBackup.pdf}}{10108\_InfrastructureBackup.pdf}{InfrastructureBackup}{Azure Stack}

\gxs{Capacity}{\includegraphics[width=1cm,valign=t]{10109_Capacity.pdf}}{10109\_Capacity.pdf}{Capacity}{Azure Stack}

\gxs{Offers}{\includegraphics[width=1cm,valign=t]{10110_Offers.pdf}}{10110\_Offers.pdf}{Offers}{Azure Stack}

\gxs{User Subscriptions}{\includegraphics[width=1cm,valign=t]{10111_UserSubscriptions.pdf}}{10111\_UserSubscriptions.pdf}{UserSubscriptions}{Azure Stack}

\gxs{Plans}{\includegraphics[width=1cm,valign=t]{10113_Plans.pdf}}{10113\_Plans.pdf}{Plans}{Azure Stack}

\gxs{Azure Stack}{\includegraphics[width=1cm,valign=t]{10114_AzureStack.pdf}}{10114\_AzureStack.pdf}{AzrStack}{Azure Stack}

\gxs{Updates}{\includegraphics[width=1cm,valign=t]{10115_Updates.pdf}}{10115\_Updates.pdf}{Updates}{Azure Stack}

\subsection{Azure VMware Solution}
\gxs{AVS}{\includegraphics[width=1cm,valign=t]{00524_AVS.pdf}}{00524\_AVS.pdf}{AVS}{Azure VMware Solution}

\subsection{Containers}
\gxs{Kubernetes Services}{\includegraphics[width=1cm,valign=t]{10023_KubernetesServices.pdf}}{10023\_KubernetesServices.pdf}{KuberSvcs}{Containers}

\gxs{Batch Accounts}{\includegraphics[width=1cm,valign=t]{10031_BatchAccounts.pdf}}{10031\_BatchAccounts.pdf}{BatchAccnts}{Containers}

\gxs{App Services}{\includegraphics[width=1cm,valign=t]{10035_AppServices.pdf}}{10035\_AppServices.pdf}{AppSvcs}{Containers}

\gxs{Service Fabric Clusters}{\includegraphics[width=1cm,valign=t]{10036_ServiceFabricClusters.pdf}}{10036\_ServiceFabricClusters.pdf}{SvcFabricClusters}{Containers}

\gxs{Container Instances}{\includegraphics[width=1cm,valign=t]{10104_ContainerInstances.pdf}}{10104\_ContainerInstances.pdf}{ContainerInsts}{Containers}

\gxs{Container Registries}{\includegraphics[width=1cm,valign=t]{10105_ContainerRegistries.pdf}}{10105\_ContainerRegistries.pdf}{ContainerRegistries}{Containers}

\subsection{Internet of Things}
\gxs{Digital Twins}{\includegraphics[width=1cm,valign=t]{01030_DigitalTwins.pdf}}{01030\_DigitalTwins.pdf}{DigitalTwins}{Internet of Things}

\subsection{IoT}
\gxs{Event Hubs}{\includegraphics[width=1cm,valign=t]{00039_EventHubs.pdf}}{00039\_EventHubs.pdf}{EventHubs}{IoT}

\gxs{Stream Analytics Jobs}{\includegraphics[width=1cm,valign=t]{00042_StreamAnalyticsJobs.pdf}}{00042\_StreamAnalyticsJobs.pdf}{StrmAnalyticsJobs}{IoT}

\gxs{Function Apps}{\includegraphics[width=1cm,valign=t]{10029_FunctionApps.pdf}}{10029\_FunctionApps.pdf}{FuncApps}{IoT}

\gxs{Notification Hubs}{\includegraphics[width=1cm,valign=t]{10045_NotificationHubs.pdf}}{10045\_NotificationHubs.pdf}{NotifHubs}{IoT}

\gxs{Time Series Insights Environments}{\includegraphics[width=1cm,valign=t]{10181_TimeSeriesInsightsEnvironments.pdf}}{10181\_TimeSeriesInsightsEnvironments.pdf}{TSInsightsEnvironments}{IoT}

\gxs{IoT Hub}{\includegraphics[width=1cm,valign=t]{10182_IoTHub.pdf}}{10182\_IoTHub.pdf}{IoTHub}{IoT}

\gxs{IoT Central Applications}{\includegraphics[width=1cm,valign=t]{10184_IoTCentralApplications.pdf}}{10184\_IoTCentralApplications.pdf}{IoTCentralApps}{IoT}

\gxs{Azure Maps Accounts}{\includegraphics[width=1cm,valign=t]{10185_AzureMapsAccounts.pdf}}{10185\_AzureMapsAccounts.pdf}{AzrMapsAccnts}{IoT}

\gxs{Time Series Insights Event Sources}{\includegraphics[width=1cm,valign=t]{10188_TimeSeriesInsightsEventSources.pdf}}{10188\_TimeSeriesInsightsEventSources.pdf}{TSInsightsEventSrcs}{IoT}

\gxs{Logic Apps}{\includegraphics[width=1cm,valign=t]{10201_LogicApps.pdf}}{10201\_LogicApps.pdf}{LogicApps}{IoT}

\gxs{Device Provisioning Services}{\includegraphics[width=1cm,valign=t]{10369_DeviceProvisioningServices.pdf}}{10369\_DeviceProvisioningServices.pdf}{DevProvisioningSvcs}{IoT}

\subsection{Preview}
\gxs{Azure Cloud Shell}{\includegraphics[width=1cm,valign=t]{00559_AzureCloudShell.pdf}}{00559\_AzureCloudShell.pdf}{AzrCloudShell}{Preview}

\gxs{Static Apps}{\includegraphics[width=1cm,valign=t]{01007_StaticApps.pdf}}{01007\_StaticApps.pdf}{StaticApps}{Preview}

\gxs{Azure Workbooks}{\includegraphics[width=1cm,valign=t]{02189_AzureWorkbooks.pdf}}{02189\_AzureWorkbooks.pdf}{AzrWorkbooks}{Preview}

\gxs{Private Link Hub}{\includegraphics[width=1cm,valign=t]{02209_PrivateLinkHub.pdf}}{02209\_PrivateLinkHub.pdf}{PrivateLinkHub}{Preview}

\gxs{IoT Edge}{\includegraphics[width=1cm,valign=t]{10186_IoTEdge.pdf}}{10186\_IoTEdge.pdf}{IoTEdge}{Preview}

\gxs{Azure Sphere}{\includegraphics[width=1cm,valign=t]{10190_AzureSphere.pdf}}{10190\_AzureSphere.pdf}{AzrSphere}{Preview}

\gxs{Time Series Data Sets}{\includegraphics[width=1cm,valign=t]{10198_TimeSeriesDataSets.pdf}}{10198\_TimeSeriesDataSets.pdf}{TSDataSets}{Preview}

\gxs{RTOS}{\includegraphics[width=1cm,valign=t]{10778_RTOS.pdf}}{10778\_RTOS.pdf}{RTOS}{Preview}

\gxs{Web Environment}{\includegraphics[width=1cm,valign=t]{10846_WebEnvironment.pdf}}{10846\_WebEnvironment.pdf}{WebEnvironment}{Preview}

\subsection{Networking}
\gxs{CDN Profiles}{\includegraphics[width=1cm,valign=t]{00056_CDNProfiles.pdf}}{00056\_CDNProfiles.pdf}{CDNProfs}{Networking}

\gxs{Azure Firewall Manager}{\includegraphics[width=1cm,valign=t]{00271_AzureFirewallManager.pdf}}{00271\_AzureFirewallManager.pdf}{AzrFirewallMngr}{Networking}

\gxs{Private Link}{\includegraphics[width=1cm,valign=t]{00427_PrivateLink.pdf}}{00427\_PrivateLink.pdf}{PrivateLink}{Networking}

\gxs{IP Groups}{\includegraphics[width=1cm,valign=t]{00701_IPGroups.pdf}}{00701\_IPGroups.pdf}{IPGrp}{Networking}

\gxs{Private Link Service}{\includegraphics[width=1cm,valign=t]{01105_PrivateLinkService.pdf}}{01105\_PrivateLinkService.pdf}{PrivateLinkSvc}{Networking}

\gxs{Virtual Networks}{\includegraphics[width=1cm,valign=t]{10061_VirtualNetworks.pdf}}{10061\_VirtualNetworks.pdf}{VirtNetworks}{Networking}

\gxs{Load Balancers}{\includegraphics[width=1cm,valign=t]{10062_LoadBalancers.pdf}}{10062\_LoadBalancers.pdf}{LBs}{Networking}

\gxs{Virtual Network Gateways}{\includegraphics[width=1cm,valign=t]{10063_VirtualNetworkGateways.pdf}}{10063\_VirtualNetworkGateways.pdf}{VirtNetworkGateways}{Networking}

\gxs{DNS Zones}{\includegraphics[width=1cm,valign=t]{10064_DNSZones.pdf}}{10064\_DNSZones.pdf}{DNSZones}{Networking}

\gxs{Traffic Manager Profiles}{\includegraphics[width=1cm,valign=t]{10065_TrafficManagerProfiles.pdf}}{10065\_TrafficManagerProfiles.pdf}{TrafficMngrProfs}{Networking}

\gxs{Network Watcher}{\includegraphics[width=1cm,valign=t]{10066_NetworkWatcher.pdf}}{10066\_NetworkWatcher.pdf}{NetworkWatcher}{Networking}

\gxs{Network Security Groups}{\includegraphics[width=1cm,valign=t]{10067_NetworkSecurityGroups.pdf}}{10067\_NetworkSecurityGroups.pdf}{NetworkSecGrp}{Networking}

\gxs{Public IP Addresses Classic}{\includegraphics[width=1cm,valign=t]{10068_PublicIPAddressesClassic.pdf}}{10068\_PublicIPAddressesClassic.pdf}{PublicIPAddressesClassic}{Networking}

\gxs{Public IP Addresses}{\includegraphics[width=1cm,valign=t]{10069_PublicIPAddresses.pdf}}{10069\_PublicIPAddresses.pdf}{PublicIPAddresses}{Networking}

\gxs{Route Filters}{\includegraphics[width=1cm,valign=t]{10071_RouteFilters.pdf}}{10071\_RouteFilters.pdf}{RouteFilters}{Networking}

\gxs{DDoS Protection Plans}{\includegraphics[width=1cm,valign=t]{10072_DDoSProtectionPlans.pdf}}{10072\_DDoSProtectionPlans.pdf}{DDoSProtectionPlans}{Networking}

\gxs{Front Doors}{\includegraphics[width=1cm,valign=t]{10073_FrontDoors.pdf}}{10073\_FrontDoors.pdf}{FrontDoors}{Networking}

\gxs{Virtual Networks Classic}{\includegraphics[width=1cm,valign=t]{10075_VirtualNetworksClassic.pdf}}{10075\_VirtualNetworksClassic.pdf}{VirtNetworksClassic}{Networking}

\gxs{Application Gateways}{\includegraphics[width=1cm,valign=t]{10076_ApplicationGateways.pdf}}{10076\_ApplicationGateways.pdf}{AppGateways}{Networking}

\gxs{ExpressRoute Circuits}{\includegraphics[width=1cm,valign=t]{10079_ExpressRouteCircuits.pdf}}{10079\_ExpressRouteCircuits.pdf}{ExpressRouteCircuits}{Networking}

\gxs{Network Interfaces}{\includegraphics[width=1cm,valign=t]{10080_NetworkInterfaces.pdf}}{10080\_NetworkInterfaces.pdf}{NetworkInterfaces}{Networking}

\gxs{Connections}{\includegraphics[width=1cm,valign=t]{10081_Connections.pdf}}{10081\_Connections.pdf}{Connections}{Networking}

\gxs{Route Tables}{\includegraphics[width=1cm,valign=t]{10082_RouteTables.pdf}}{10082\_RouteTables.pdf}{RouteTables}{Networking}

\gxs{Firewalls}{\includegraphics[width=1cm,valign=t]{10084_Firewalls.pdf}}{10084\_Firewalls.pdf}{Firewalls}{Networking}

\gxs{Service Endpoint Policies}{\includegraphics[width=1cm,valign=t]{10085_ServiceEndpointPolicies.pdf}}{10085\_ServiceEndpointPolicies.pdf}{SvcEndpointPolicies}{Networking}

\gxs{NAT}{\includegraphics[width=1cm,valign=t]{10310_NAT.pdf}}{10310\_NAT.pdf}{NAT}{Networking}

\gxs{Virtual WANs}{\includegraphics[width=1cm,valign=t]{10353_VirtualWANs.pdf}}{10353\_VirtualWANs.pdf}{VirtWANs}{Networking}

\gxs{Web Application Firewall PoliciesWAF}{\includegraphics[width=1cm,valign=t]{10362_WebApplicationFirewallPoliciesWAF.pdf}}{10362\_WebApplicationFirewallPoliciesWAF.pdf}{WebAppFirewallPoliciesWAF}{Networking}

\gxs{Proximity Placement Groups}{\includegraphics[width=1cm,valign=t]{10365_ProximityPlacementGroups.pdf}}{10365\_ProximityPlacementGroups.pdf}{ProximityPlacementGrp}{Networking}

\gxs{Reserved IP Addresses Classic}{\includegraphics[width=1cm,valign=t]{10371_ReservedIPAddressesClassic.pdf}}{10371\_ReservedIPAddressesClassic.pdf}{ReservedIPAddressesClassic}{Networking}

\gxs{Public IP Prefixes}{\includegraphics[width=1cm,valign=t]{10372_PublicIPPrefixes.pdf}}{10372\_PublicIPPrefixes.pdf}{PublicIPPrefixes}{Networking}

\subsection{Other}
\gxs{Windows Virtual Desktop}{\includegraphics[width=1cm,valign=t]{00327_WindowsVirtualDesktop.pdf}}{00327\_WindowsVirtualDesktop.pdf}{WinVirtDesktop}{Other}

\gxs{Detonation}{\includegraphics[width=1cm,valign=t]{00378_Detonation.pdf}}{00378\_Detonation.pdf}{Detonation}{Other}

\gxs{Internet Analyzer Profiles}{\includegraphics[width=1cm,valign=t]{00469_InternetAnalyzerProfiles.pdf}}{00469\_InternetAnalyzerProfiles.pdf}{INetAnalyzerProfs}{Other}

\gxs{Universal Print}{\includegraphics[width=1cm,valign=t]{00571_UniversalPrint.pdf}}{00571\_UniversalPrint.pdf}{UniverPrint}{Other}

\gxs{Peering Service}{\includegraphics[width=1cm,valign=t]{00970_PeeringService.pdf}}{00970\_PeeringService.pdf}{PeeringSvc}{Other}

\gxs{Instance Pools}{\includegraphics[width=1cm,valign=t]{10139_InstancePools.pdf}}{10139\_InstancePools.pdf}{InstPools}{Other}

\subsection{Storage}
\gxs{Recovery Services Vaults}{\includegraphics[width=1cm,valign=t]{00017_RecoveryServicesVaults.pdf}}{00017\_RecoveryServicesVaults.pdf}{RecovSvcsVlts}{Storage}

\gxs{Azure Stack Edge}{\includegraphics[width=1cm,valign=t]{00691_AzureStackEdge.pdf}}{00691\_AzureStackEdge.pdf}{AzrStackEdge}{Storage}

\gxs{Azure HCP Cache}{\includegraphics[width=1cm,valign=t]{00776_AzureHCPCache.pdf}}{00776\_AzureHCPCache.pdf}{AzrHCPCache}{Storage}

\gxs{Storage Accounts}{\includegraphics[width=1cm,valign=t]{10086_StorageAccounts.pdf}}{10086\_StorageAccounts.pdf}{StorAccnts}{Storage}

\gxs{Storage Accounts Classic}{\includegraphics[width=1cm,valign=t]{10087_StorageAccountsClassic.pdf}}{10087\_StorageAccountsClassic.pdf}{StorAccntsClassic}{Storage}

\gxs{StorSimple Device Managers}{\includegraphics[width=1cm,valign=t]{10089_StorSimpleDeviceManagers.pdf}}{10089\_StorSimpleDeviceManagers.pdf}{StorSimpleDevMngrs}{Storage}

\gxs{Data Lake Storage Gen1}{\includegraphics[width=1cm,valign=t]{10090_DataLakeStorageGen1.pdf}}{10090\_DataLakeStorageGen1.pdf}{DataLakeStorGenOne}{Storage}

\gxs{StorSimple Data Managers}{\includegraphics[width=1cm,valign=t]{10092_StorSimpleDataManagers.pdf}}{10092\_StorSimpleDataManagers.pdf}{StorSimpleDataMngrs}{Storage}

\gxs{Storage Sync Services}{\includegraphics[width=1cm,valign=t]{10093_StorageSyncServices.pdf}}{10093\_StorageSyncServices.pdf}{StorSyncSvcs}{Storage}

\gxs{Data Box}{\includegraphics[width=1cm,valign=t]{10094_DataBox.pdf}}{10094\_DataBox.pdf}{DataBox}{Storage}

\gxs{Data Box Edge}{\includegraphics[width=1cm,valign=t]{10095_DataBoxEdge.pdf}}{10095\_DataBoxEdge.pdf}{DataBoxEdge}{Storage}

\gxs{Azure NetApp Files}{\includegraphics[width=1cm,valign=t]{10096_AzureNetAppFiles.pdf}}{10096\_AzureNetAppFiles.pdf}{AzrNetAppFiles}{Storage}

\gxs{Data Share Invitations}{\includegraphics[width=1cm,valign=t]{10097_DataShareInvitations.pdf}}{10097\_DataShareInvitations.pdf}{DataShareInvitations}{Storage}

\gxs{Data Shares}{\includegraphics[width=1cm,valign=t]{10098_DataShares.pdf}}{10098\_DataShares.pdf}{DataShares}{Storage}

\gxs{Import Export Jobs}{\includegraphics[width=1cm,valign=t]{10100_ImportExportJobs.pdf}}{10100\_ImportExportJobs.pdf}{ImportExportJobs}{Storage}

\subsection{AI + Machine Learning}
\gxs{Cognitive Services}{\includegraphics[width=1cm,valign=t]{10162_CognitiveServices.pdf}}{10162\_CognitiveServices.pdf}{CognitiveSvcs}{AI + Machine Learning}

\gxs{Bot Services}{\includegraphics[width=1cm,valign=t]{10165_BotServices.pdf}}{10165\_BotServices.pdf}{BotSvcs}{AI + Machine Learning}

\subsection{App Services}
\gxs{App Service Plans}{\includegraphics[width=1cm,valign=t]{00046_AppServicePlans.pdf}}{00046\_AppServicePlans.pdf}{AppSvcPlans}{App Services}

\gxs{App Service Certificates}{\includegraphics[width=1cm,valign=t]{00049_AppServiceCertificates.pdf}}{00049\_AppServiceCertificates.pdf}{AppSvcCerts}{App Services}

\gxs{App Service Domains}{\includegraphics[width=1cm,valign=t]{00050_AppServiceDomains.pdf}}{00050\_AppServiceDomains.pdf}{AppSvcDomains}{App Services}

\gxs{CDN Profiles}{\includegraphics[width=1cm,valign=t]{00056_CDNProfiles.pdf}}{00056\_CDNProfiles.pdf}{CDNProfs}{App Services}

\gxs{App Services}{\includegraphics[width=1cm,valign=t]{10035_AppServices.pdf}}{10035\_AppServices.pdf}{AppSvcs}{App Services}

\gxs{API Management Services}{\includegraphics[width=1cm,valign=t]{10042_APIManagementServices.pdf}}{10042\_APIManagementServices.pdf}{APIMngmtSvcs}{App Services}

\gxs{Search Services}{\includegraphics[width=1cm,valign=t]{10044_SearchServices.pdf}}{10044\_SearchServices.pdf}{SrchSvcs}{App Services}

\gxs{Notification Hubs}{\includegraphics[width=1cm,valign=t]{10045_NotificationHubs.pdf}}{10045\_NotificationHubs.pdf}{NotifHubs}{App Services}

\gxs{App Service Environments}{\includegraphics[width=1cm,valign=t]{10047_AppServiceEnvironments.pdf}}{10047\_AppServiceEnvironments.pdf}{AppSvcEnvironments}{App Services}

\subsection{Databases}
\gxs{SQL Data Warehouses}{\includegraphics[width=1cm,valign=t]{00036_SQLDataWarehouses.pdf}}{00036\_SQLDataWarehouses.pdf}{SQLDataWarehouses}{Databases}

\gxs{Azure Synapse Analytics}{\includegraphics[width=1cm,valign=t]{00606_AzureSynapseAnalytics.pdf}}{00606\_AzureSynapseAnalytics.pdf}{AzrSynapseAnalytics}{Databases}

\gxs{Azure SQL}{\includegraphics[width=1cm,valign=t]{02390_AzureSQL.pdf}}{02390\_AzureSQL.pdf}{AzrSQL}{Databases}

\gxs{SSIS Lift And Shift IR}{\includegraphics[width=1cm,valign=t]{02392_SSISLiftAndShiftIR.pdf}}{02392\_SSISLiftAndShiftIR.pdf}{SSISLiftAndShiftIR}{Databases}

\gxs{Azure Cosmos DB}{\includegraphics[width=1cm,valign=t]{10121_AzureCosmosDB.pdf}}{10121\_AzureCosmosDB.pdf}{AzrCosmosDB}{Databases}

\gxs{Azure Database MySQL Server}{\includegraphics[width=1cm,valign=t]{10122_AzureDatabaseMySQLServer.pdf}}{10122\_AzureDatabaseMySQLServer.pdf}{AzrDBMySQLServer}{Databases}

\gxs{Azure Database MariaDB Server}{\includegraphics[width=1cm,valign=t]{10123_AzureDatabaseMariaDBServer.pdf}}{10123\_AzureDatabaseMariaDBServer.pdf}{AzrDBMariaDBServer}{Databases}

\gxs{Azure SQL VM}{\includegraphics[width=1cm,valign=t]{10124_AzureSQLVM.pdf}}{10124\_AzureSQLVM.pdf}{AzrSQLVM}{Databases}

\gxs{Data Factory}{\includegraphics[width=1cm,valign=t]{10126_DataFactory.pdf}}{10126\_DataFactory.pdf}{DataFactory}{Databases}

\gxs{Virtual Clusters}{\includegraphics[width=1cm,valign=t]{10127_VirtualClusters.pdf}}{10127\_VirtualClusters.pdf}{VirtClusters}{Databases}

\gxs{Elastic Job Agents}{\includegraphics[width=1cm,valign=t]{10128_ElasticJobAgents.pdf}}{10128\_ElasticJobAgents.pdf}{ElasticJobAgents}{Databases}

\gxs{SQL Database}{\includegraphics[width=1cm,valign=t]{10130_SQLDatabase.pdf}}{10130\_SQLDatabase.pdf}{SQLDB}{Databases}

\gxs{Azure Database PostgreSQL Server}{\includegraphics[width=1cm,valign=t]{10131_AzureDatabasePostgreSQLServer.pdf}}{10131\_AzureDatabasePostgreSQLServer.pdf}{AzrDBPostgreSQLServer}{Databases}

\gxs{SQL Server}{\includegraphics[width=1cm,valign=t]{10132_SQLServer.pdf}}{10132\_SQLServer.pdf}{SQLServer}{Databases}

\gxs{Azure Database Migration Services}{\includegraphics[width=1cm,valign=t]{10133_AzureDatabaseMigrationServices.pdf}}{10133\_AzureDatabaseMigrationServices.pdf}{AzrDBMigratSvcs}{Databases}

\gxs{SQL Elastic Pools}{\includegraphics[width=1cm,valign=t]{10134_SQLElasticPools.pdf}}{10134\_SQLElasticPools.pdf}{SQLElasticPools}{Databases}

\gxs{Managed Database}{\includegraphics[width=1cm,valign=t]{10135_ManagedDatabase.pdf}}{10135\_ManagedDatabase.pdf}{MngdDB}{Databases}

\gxs{SQL Managed Instance}{\includegraphics[width=1cm,valign=t]{10136_SQLManagedInstance.pdf}}{10136\_SQLManagedInstance.pdf}{SQLMngdInst}{Databases}

\gxs{Azure SQL Server Stretch Databases}{\includegraphics[width=1cm,valign=t]{10137_AzureSQLServerStretchDatabases.pdf}}{10137\_AzureSQLServerStretchDatabases.pdf}{AzrSQLServerStretchDBs}{Databases}

\gxs{Cache Redis}{\includegraphics[width=1cm,valign=t]{10137_CacheRedis.pdf}}{10137\_CacheRedis.pdf}{CacheRedis}{Databases}

\gxs{Instance Pools}{\includegraphics[width=1cm,valign=t]{10139_InstancePools.pdf}}{10139\_InstancePools.pdf}{InstPools}{Databases}

\gxs{Azure Data Explorer Clusters}{\includegraphics[width=1cm,valign=t]{10145_AzureDataExplorerClusters.pdf}}{10145\_AzureDataExplorerClusters.pdf}{AzrDataExplorerClusters}{Databases}

\subsection{Identity}
\gxs{Azure Active Directory}{\includegraphics[width=1cm,valign=t]{10221_AzureActiveDirectory.pdf}}{10221\_AzureActiveDirectory.pdf}{AzrActiveDir}{Identity}

\gxs{Azure AD Domain Services}{\includegraphics[width=1cm,valign=t]{10222_AzureADDomainServices.pdf}}{10222\_AzureADDomainServices.pdf}{AzrADDomainSvcs}{Identity}

\gxs{Groups}{\includegraphics[width=1cm,valign=t]{10223_Groups.pdf}}{10223\_Groups.pdf}{Grp}{Identity}

\gxs{Active Directory Connect Health}{\includegraphics[width=1cm,valign=t]{10224_ActiveDirectoryConnectHealth.pdf}}{10224\_ActiveDirectoryConnectHealth.pdf}{ActiveDirConnectHealth}{Identity}

\gxs{Enterprise Applications}{\includegraphics[width=1cm,valign=t]{10225_EnterpriseApplications.pdf}}{10225\_EnterpriseApplications.pdf}{EntrprsApps}{Identity}

\gxs{Azure AD B2C}{\includegraphics[width=1cm,valign=t]{10228_AzureADB2C.pdf}}{10228\_AzureADB2C.pdf}{AzrADBTwoC}{Identity}

\gxs{Users}{\includegraphics[width=1cm,valign=t]{10230_Users.pdf}}{10230\_Users.pdf}{Users}{Identity}

\gxs{App Registrations}{\includegraphics[width=1cm,valign=t]{10232_AppRegistrations.pdf}}{10232\_AppRegistrations.pdf}{AppRegistrations}{Identity}

\gxs{Identity Governance}{\includegraphics[width=1cm,valign=t]{10235_IdentityGovernance.pdf}}{10235\_IdentityGovernance.pdf}{IdentGovern}{Identity}

\subsection{Migrate}
\gxs{Recovery Services Vaults}{\includegraphics[width=1cm,valign=t]{00017_RecoveryServicesVaults.pdf}}{00017\_RecoveryServicesVaults.pdf}{RecovSvcsVlts}{Migrate}

\gxs{Data Box}{\includegraphics[width=1cm,valign=t]{10094_DataBox.pdf}}{10094\_DataBox.pdf}{DataBox}{Migrate}

\gxs{Data Box Edge}{\includegraphics[width=1cm,valign=t]{10095_DataBoxEdge.pdf}}{10095\_DataBoxEdge.pdf}{DataBoxEdge}{Migrate}

\gxs{Azure Migrate}{\includegraphics[width=1cm,valign=t]{10281_AzureMigrate.pdf}}{10281\_AzureMigrate.pdf}{AzrMigrate}{Migrate}


\end{multicols}\clearpage

%%%%%%%%%%%%% INDEXES
\normalsize

\appendix
\section{Mapping of icon names to macros names}
In order to have reasonably short macro names, the following mappings have been made by \texttt{aws2tex}, being careful not to introduce any naming collisions.
In addition, \LaTeX\ macros cannot contain numbers, so all numbers have been mapped to words, i.e. 1 $\rightarrow$ One, 5 $\rightarrow$ Five, etc.
The below mappings are applied in order, followed by the number mappings.
{\tiny
\begin{multicols}{3}
  \begin{enumerate}
  \item example
  \end{enumerate}
\end{multicols}
}

\printindex
\printindex[macros]

\end{document}
